\begin{dang}{Máy phát điện xoay chiều một pha và đoạn mạch RLC}
\Anlythuyet{	\Binhluan{
		\begin{itemize}
		\item Cường độ hiệu dụng trong mạch là: $I =$ $\dfrac{E}{{\sqrt {{R^2} + {{\left( {{Z_L} - {Z_C}} \right)}^2}} }}$
		\item Với $\left\{ \begin{array}{l}
		f =np \Rightarrow \omega  = 2\pi f \Rightarrow {Z_L} = \omega L,\,{Z_C} = \dfrac{1}{{\omega C}}\\[0.3cm]
		E = \dfrac{{N.2\pi f.{\Phi _0}}}{{\sqrt 2 }}
		\end{array} \right.$
		\item Khi $n’ = kn$ thì  $E' =kE,\ Z{'_L} = kZ_L,\ Z'_C = \dfrac{Z_C}{k}\\
		 \Rightarrow I' =\dfrac{{kE}}{{\sqrt {{R^2} + {{\left( {kZ_L - \dfrac{Z_C}{k}} \right)}^2}} }} \Rightarrow \dfrac{{I'}}{I} =k\dfrac{{\sqrt {R^2 +{{\left(Z_L - Z_C \right)}^2}}}}{{\sqrt {{R^2} + {{\left( {kZ_L - \dfrac{Z_C}{k}}\right)}^2}}}}$
	\end{itemize}
}
{	Với những bài toán thay đổi tốc độ quay rôto qua ba giá trị ta nên lập bảng chuẩn hoá số liệu}}
\end{dang}
